\section*{Chapter \ref{ch:g10:freshwater-resources} --- Freshwater Resources}

\begin{solution}{exo:g10:freshwater-resources:1}
A \omterm{def:g10:freshwater-resources:stock}{stock} is the amount of water held in a reservoir at a given time, in
\unit{km^3} (or kilograms). A \omterm{def:g10:freshwater-resources:stock}{flux} is the amount that passes from one
reservoir to another per unit time, in \unit{km^3} per year. The \omterm{def:g10:freshwater-resources:residence-time}{residence
time} is the average time water spends in a reservoir, in days or years; in
steady state it is the \omterm{def:g10:freshwater-resources:stock}{stock} divided by the \omterm{def:g10:freshwater-resources:stock}{flux}.
\end{solution}

\begin{solution}{exo:g10:freshwater-resources:2}
About \qty{2.5}{\percent} of all water is fresh. Of the fresh water,
\qty{68.7}{\percent} is frozen in ice caps, glaciers and permanent snow, and
\qty{30.1}{\percent} is underground. Rivers hold the least,
\num{2120}~\unit{km^3}.
\end{solution}

\begin{solution}{exo:g10:freshwater-resources:3}
\omterm{def:g10:freshwater-resources:renewable}{Renewable fresh water} is the fresh water renewed each year by precipitation
and flowing through a region as runoff and groundwater: a volume per year,
a \omterm{def:g10:freshwater-resources:stock}{flux}. \omterm{def:g10:freshwater-resources:fossil-groundwater}{Fossil groundwater} is water that entered an aquifer long ago and that
today's recharge does not renew: it is a store that can only be drawn down,
a \omterm{def:g10:freshwater-resources:stock}{stock}.
\end{solution}

\begin{solution}{exo:g10:freshwater-resources:4}
$91\,000/2120 \approx 43$ times more.
\end{solution}

\begin{solution}{exo:g10:freshwater-resources:5}
The atmosphere renews its water fastest (about \num{9.5} days), the ocean
slowest among the computed reservoirs (about \num{3200} years); deep
groundwater is older still. $3200 \times 365.25/9.5 \approx \num{120000}$
times longer.
\end{solution}

\begin{solution}{exo:g10:freshwater-resources:6}
In: $413\,000 + 73\,000 = \num{486000}$~\unit{km^3} a year, so
$T = 12\,700/486\,000 \approx \num{0.0261}$ year $\approx \num{9.5}$ days.
Out: $373\,000 + 113\,000 = \num{486000}$~\unit{km^3}, the same: the air is in
steady state, and the \omterm{def:g10:freshwater-resources:residence-time}{residence time} is the same.
\end{solution}

\begin{solution}{exo:g10:freshwater-resources:7}
$1\,335\,040/413 \approx \num{3230}$ years. In steady state the ocean
receives as much as it loses: $373 + 40 = 413$ thousand~\unit{km^3} a year,
so the \omterm{def:g10:freshwater-resources:residence-time}{residence time} computed with the inputs is the same.
\end{solution}

\begin{solution}{exo:g10:freshwater-resources:8}
The land receives \num{113000}~\unit{km^3} a year; \num{73000} return to the
air by evaporation and transpiration and \num{40000} run off to the sea:
$73\,000 + 40\,000 = \num{113000}$. The share returned to the air is
$73/113 \approx \qty{65}{\percent}$.
\end{solution}

\begin{solution}{exo:g10:freshwater-resources:9}
$273.2 \times 10^6 \times 1233.5 \approx \qty{3.37e11}{m^3}$, that is
\qty{337}{km^3}: almost as much as all the water the Aral Sea still held in
1987.
\end{solution}

\begin{solution}{exo:g10:freshwater-resources:10}
$15.8 \times 0.3048 \approx \qty{4.82}{m}$. Over the 65 years from 1950 to
2015, $482/65 \approx \qty{7.4}{cm}$ a year.
\end{solution}

\begin{solution}{exo:g10:freshwater-resources:11}
$2/0.01 = 200$ years. The \omterm{def:g10:freshwater-resources:stock}{stock} now falls by $0.05 - 0.01 =
\qty{0.04}{km^3}$ a year, so it lasts $2/0.04 = 50$ years.
\end{solution}

\begin{solution}{exo:g10:freshwater-resources:12}
1960: $1093/71.1 \approx 15$ years; 2010: $98.1/11.4 \approx 8.6$ years.
Evaporation removes about the same thickness of water each year from every
square kilometre of the lake; as the lake shrank, its area shrank, and so did
the total volume evaporated.
\end{solution}

\begin{solution}{exo:g10:freshwater-resources:13}
(1) The net loss is $1.2 - 0.2 = \qty{1.0}{m}$ a year over
$100 \times 10^6$~\unit{m^2}: $10^8$~\unit{m^3}, \qty{0.1}{km^3} a year must
flow in. (2) The inflow drops to \qty{0.05}{km^3}: the lake loses
$\qty{0.05}{km^3} = 5 \times 10^7$~\unit{m^3}, a layer of
$5 \times 10^7/10^8 = \qty{0.5}{m}$. (3) As the level falls, the shores
come out of the water and the lake's area shrinks: less surface, less
evaporation, until evaporation and inflow balance again.
\end{solution}

\begin{solution}{exo:g10:freshwater-resources:14}
The aquifer holds $100 \times 0.20 = \qty{20}{m}$ of water, that is
\qty{20000}{mm}: at \qty{1}{mm} a year, \num{20000} years. The deep water of
the Sahara aquifer is \num{200000} to a million years old, ten to fifty times
more: it was recharged during wetter past climates.
\end{solution}

\begin{solution}{exo:g10:freshwater-resources:15}
$337/453\,000 \approx \qty{0.00074}{km}$, about \qty{0.74}{m} of water. The
water level fell \qty{4.8}{m}, about six and a half times more. The water
fills only the pores of the rock (and not all of it drains out): when the
water table falls by a metre, the water released is a layer much thinner than
a metre, here about $0.74/4.8 \approx \num{0.15}$~m.
\end{solution}

\begin{solution}{pb:g10:freshwater-resources:1}
\textbf{1.} $1093/68\,000 \approx \qty{0.016}{km}$, about \qty{16}{m}.

\textbf{2.} $1093 \times 10^9 \approx \qty{1.09e12}{m^3}$.

\textbf{3.} \qty{10}{g/L} is \qty{10}{kg/m^3}: $1.093 \times 10^{12} \times 10
\approx \qty{1.1e13}{kg}$, about \num{11} billion tonnes.

\textbf{4.} $1093/71.1 \approx 15$ years.

\textbf{5.} No river leaves it: its water can only leave as vapour (or by
seeping into the ground), while the salt stays.

\textbf{6.} $71.1/68\,000 \approx \qty{0.00105}{km}$, about \qty{1.05}{m}.

\textbf{7.} $9.4/68\,000 \approx \qty{0.00014}{km}$, about \qty{0.14}{m}.

\textbf{8.} $71.1 - 9.4 = \qty{61.7}{km^3}$ a year.

\textbf{9.} $79 + 37 = \qty{116}{km^3}$: $61.7/116 \approx \qty{53}{\percent}$,
about half.

\textbf{10.} Inputs $1.0 + 6.2 = \qty{7.2}{km^3}$; output \qty{36.8}{km^3};
net change $7.2 - 36.8 = \qty{-29.6}{km^3}$.

\textbf{11.} $(502.7 - 327.2)/5 = \qty{35.1}{km^3}$ a year: the same order as
the \qty{29.6}{km^3} of the 1987 budget (single years differ: evaporation was
\qty{47.9}{km^3} in 1984).

\textbf{12.} $36.8/0.00105 \approx \qty{35000}{km^2}$.

\textbf{13.} $35\,000/68\,000 \approx \qty{52}{\percent}$, about half.

\textbf{14.} The lake had shrunk: half the surface evaporates about half the
water.

\textbf{15.} $345.6/1093 \approx \qty{32}{\percent}$, about a third.

\textbf{16.} $10 \times 1093/345.6 \approx \qty{31.6}{g/L}$.

\textbf{17.} The measured salinity is lower: part of the salt had left the
water, crystallised on the exposed bed and carried off by the wind as salty
dust.

\textbf{18.} About \qty{1.05}{m} a year, the thickness evaporated.

\textbf{19.} $7.2/35\,000 \approx \qty{0.0002}{km}$, about \qty{0.20}{m}.

\textbf{20.} $1.05 - 0.20 \approx 0.85$ (\num{0.84} with the unrounded values):
\textbf{in 1987 the level of the Aral Sea fell by about \qty{0.8}{m}, close to
a metre in a single year}.
\end{solution}
